% =====================================================================
\chapter{Soluciones de los ejercicios}
\label{ap:soluciones}
% =====================================================================

Cada solución se presenta en cinco partes: datos, planteamiento, desarrollo, resultado y comprobación. En los pasos intermedios se redondea a cuatro cifras significativas salvo que se indique otra cosa; las raíces exactas ($\sqrt{2}$, $\sqrt{5}$) se conservan hasta el final cuando es posible. En las demostraciones, la comprobación consiste en un caso numérico o en un contraejemplo que muestra que las hipótesis son necesarias.

% ---------------------------------------------------------------------
\section{Tema 1}

\begin{solucion}{eje:1-vf}
\planteamiento aplicar el \cref{teo:dimension}, el \cref{lem:steinitz} y el \cref{teo:filas-columnas}.
\pasos
\begin{enumerate}[label=(\alph*)]
  \item \textbf{Falsa.} $\R^{2}=\gen\{e_1,e_2\}$ y, por el \cref{lem:steinitz}, tres vectores de $\R^{2}$ son dependientes. Por ejemplo, $(1,1)=(1,0)+(0,1)$.
  \item \textbf{Verdadera.} Las $n$ columnas son vectores independientes de $\R^{m}$, generado por $m$ vectores; si fuera $n>m$, el \cref{lem:steinitz} las haría dependientes.
  \item \textbf{Falsa.} El rango es el número de columnas independientes, como mucho~3, y también el de filas independientes (\cref{teo:filas-columnas}). Luego $\rango(A)\le\min\{5,3\}=3<4$.
\end{enumerate}
\resultado \textbf{(a) falsa; (b) verdadera; (c) falsa.}
\comprobacion en (c), una matriz de $5\times 3$ con columnas $e_1,e_2,e_3\in\R^{5}$ alcanza el rango máximo, que es~3.
\end{solucion}

\begin{solucion}{eje:1-rango1}
\datos $u=(1,2,3)$, $v=(1,-1)$.
\planteamiento $(uv\T)_{ij}=u_iv_j$; cada fila es $u_i\,v\T$ y cada columna es $v_j\,u$.
\pasos
\begin{enumerate}[label=\arabic*.]
  \item Fila 1: $1\cdot(1,-1)=(1,-1)$; fila 2: $2\cdot(1,-1)=(2,-2)$; fila 3: $3\cdot(1,-1)=(3,-3)$.
  \[
  uv\T=\begin{pmatrix}1&-1\\2&-2\\3&-3\end{pmatrix}.
  \]
  \item La columna~2 es $-1$ veces la columna~1, que es no nula: $\rango(uv\T)=1$ y $\Col(uv\T)=\gen\{(1,2,3)\}$.
  \item Núcleo: $(uv\T)x=u\,(v\T x)=0$ si y solo si $v\T x=x_1-x_2=0$. Luego $\Nuc(uv\T)=\gen\{(1,1)\}$.
\end{enumerate}
\resultado \textbf{$uv\T$ es de $3\times 2$, tiene rango~1; base de $\Col$: $\{(1,2,3)\}$; base de $\Nuc$: $\{(1,1)\}$.}
\comprobacion rango y nulidad: $1+1=2$, el número de columnas (\cref{teo:rango-nulidad}). Además, $uv\T(1,1)=(1-1,\,2-2,\,3-3)=0$.
\end{solucion}

\begin{solucion}{eje:1-mml}
\datos $x_1=(2,-1,3)$, $x_2=(1,1,-2)$, $x_3=(3,-3,8)$; $y=(1,-2,5)$, $z_1=(1,1,1)$, $z_2=(1,2,3)$, $z_3=(2,-1,1)$.
\planteamiento (a) buscar $a,b$ con $x_3=ax_1+bx_2$, o probar que no existen; (b) resolver el sistema $c_1z_1+c_2z_2+c_3z_3=y$.
\pasos
\begin{enumerate}[label=\arabic*.]
  \item (a) Primeras dos componentes de $x_3=ax_1+bx_2$: $2a+b=3$ y $-a+b=-3$. Restando, $3a=6$, luego $a=2$ y $b=3-4=-1$.
  \item Tercera componente: $3a-2b=6+2=8$, que coincide con la de $x_3$. Por tanto, $x_3=2x_1-x_2$: \textbf{son dependientes}, con relación $2x_1-x_2-x_3=0$.
  \item (b) Sistema: $c_1+c_2+2c_3=1$, $c_1+2c_2-c_3=-2$, $c_1+3c_2+c_3=5$.
  \item Restando la primera ecuación a la segunda y a la tercera: $c_2-3c_3=-3$ y $2c_2-c_3=4$.
  \item De la primera de ellas, $c_2=3c_3-3$; sustituyendo, $6c_3-6-c_3=4$, luego $c_3=2$ y $c_2=3$.
  \item $c_1=1-c_2-2c_3=1-3-4=-6$.
  \item Unicidad: la matriz cuyas columnas son $z_1,z_2,z_3$ es $Z=\begin{pmatrix}1&1&2\\1&2&-1\\1&3&1\end{pmatrix}$. Desarrollando por la primera fila: $1\cdot(2\cdot 1-(-1)\cdot 3)-1\cdot(1\cdot 1-(-1)\cdot 1)+2\cdot(1\cdot 3-2\cdot 1)=5-2+2=5\ne 0$. Las $z_i$ son una base de $\R^{3}$ y las coordenadas son únicas (\cref{prop:coordenadas}).
\end{enumerate}
\resultado \textbf{(a) dependientes: $x_3=2x_1-x_2$. (b) $y=-6z_1+3z_2+2z_3$, de forma única porque $\det Z=5\ne 0$.}
\comprobacion $-6(1,1,1)+3(1,2,3)+2(2,-1,1)=(-6+3+4,\,-6+6-2,\,-6+9+2)=(1,-2,5)$.
\end{solucion}

\begin{solucion}{eje:1-u4}
\datos $v_1=(1,-1,0)$, $v_2=(1,0,-1)$, $v_3=(0,1,-1)$, $v_4=(1,1,-2)$.
\planteamiento buscar relaciones de dependencia; el número de independientes es el rango de la matriz de $3\times 4$ con esas columnas.
\pasos
\begin{enumerate}[label=\arabic*.]
  \item $v_2-v_1=(1-1,\,0+1,\,-1-0)=(0,1,-1)=v_3$. Relación: $v_1-v_2+v_3=0$.
  \item $v_2+v_3=(1+0,\,0+1,\,-1-1)=(1,1,-2)=v_4$. Relación: $v_2+v_3-v_4=0$.
  \item $v_1$ y $v_2$ son independientes: $c_1(1,-1,0)+c_2(1,0,-1)=0$ obliga a $c_1=0$ (segunda componente) y a $c_2=0$ (tercera componente).
  \item Por los pasos 1 y 2, $v_3$ y $v_4$ son combinación de $v_1$ y $v_2$: el rango es~2 y $\{v_1,v_2\}$ es una base.
  \item Las componentes de cada $v_i$ suman cero: $1-1+0$, $1+0-1$, $0+1-1$ y $1+1-2$. El subespacio es el plano $x+y+z=0$.
\end{enumerate}
\resultado \textbf{como máximo 2 independientes; base $\{v_1,v_2\}$; el subespacio es el plano $x+y+z=0$.}
\comprobacion por el \cref{teo:rango-nulidad}, la matriz de $3\times 4$ tiene núcleo de dimensión $4-2=2$, generado por los coeficientes de las relaciones $(1,-1,1,0)$ y $(0,1,1,-1)$. El plano tiene dimensión $3-1=2$.
\end{solucion}

\begin{solucion}{eje:1-nucleo}
\datos $A=\begin{pmatrix}1&2&0&1\\2&4&1&3\\3&6&1&4\end{pmatrix}$.
\planteamiento eliminación gaussiana para obtener los pivotes; las columnas pivote de $A$ forman una base de $\Col(A)$ y cada variable libre da un vector de la base de $\Nuc(A)$.
\pasos
\begin{enumerate}[label=\arabic*.]
  \item $F_2\leftarrow F_2-2F_1$: $(0,0,1,1)$. $F_3\leftarrow F_3-3F_1$: $(0,0,1,1)$. $F_3\leftarrow F_3-F_2$: $(0,0,0,0)$.
  \item Forma escalonada: $\begin{pmatrix}1&2&0&1\\0&0&1&1\\0&0&0&0\end{pmatrix}$. Pivotes en las columnas~1 y~3: $\rango(A)=2$.
  \item Base de $\Col(A)$: columnas 1 y 3 de $A$, $\{(1,2,3),\,(0,1,1)\}$.
  \item Núcleo: $x_3=-x_4$ y $x_1=-2x_2-x_4$, con $x_2$ y $x_4$ libres. Con $(x_2,x_4)=(1,0)$: $(-2,1,0,0)$. Con $(x_2,x_4)=(0,1)$: $(-1,0,-1,1)$.
\end{enumerate}
\resultado \textbf{$\rango(A)=2$.} Una base de $\Col(A)$ la forman $(1,2,3)$ y $(0,1,1)$, y una base de $\Nuc(A)$, los vectores $(-2,1,0,0)$ y $(-1,0,-1,1)$. Se cumple $2+2=4$, el número de columnas.
\comprobacion $A(-2,1,0,0)=(-2+2,\,-4+4,\,-6+6)=0$ y $A(-1,0,-1,1)=(-1+0+1,\,-2-1+3,\,-3-1+4)=0$. La fila~3 de $A$ es la suma de las filas~1 y~2.
\end{solucion}

\begin{solucion}{eje:1-producto}
\datos $B\in\R^{m\times r}$, $A\in\R^{r\times n}$; $\Delta W^{\star}\in\R^{d\times d}$ con $\rango(\Delta W^{\star})>r$.
\planteamiento comparar el espacio columna de $BA$ con el de $B$, y el espacio de filas de $BA$ con el de $A$.
\pasos
\begin{enumerate}[label=\arabic*.]
  \item Para todo $x$, $(BA)x=B(Ax)\in\Col(B)$. Luego $\Col(BA)\subseteq\Col(B)$ y $\rango(BA)\le\rango(B)$.
  \item Cada fila de $BA$ es (fila de $B$)$\cdot A$, combinación lineal de las filas de $A$. Por tanto, $\Col((BA)\T)\subseteq\Col(A\T)$ y, por el \cref{teo:filas-columnas}, $\rango(BA)=\rango((BA)\T)\le\rango(A\T)=\rango(A)$.
  \item Juntando, $\rango(BA)\le\min\{\rango(A),\rango(B)\}$.
  \item $B$ tiene $r$ columnas, luego $\rango(B)\le r$ y $\rango(BA)\le r$.
  \item Si fuera $BA=\Delta W^{\star}$, ambas matrices tendrían el mismo rango; pero $\rango(BA)\le r<\rango(\Delta W^{\star})$.
\end{enumerate}
\resultado \textbf{$\rango(BA)\le\min\{\rango A,\rango B\}\le r$, y ninguna corrección LoRA de rango $r$ iguala una matriz de rango mayor que $r$.} El error mínimo que se comete se calcula en el \cref{eje:5-lora}. \hfill$\blacksquare$
\comprobacion $B=(1,2)\T$ y $A=(3,4)$ dan $BA=\begin{pmatrix}3&4\\6&8\end{pmatrix}$, de rango $1=r$.
\end{solucion}

% ---------------------------------------------------------------------
\section{Tema 2}

\begin{solucion}{eje:2-equiv}
\planteamiento identificar en cada caso una condición del \cref{teo:invertible} que se cumpla o falle.
\pasos
\begin{enumerate}[label=(\alph*)]
  \item $M$ no es invertible: $\rango(M)=2<3$ y $M(1,1,-1)=0$ (\cref{ej:M}).
  \item $Q_\theta$ es invertible: $\det(Q_\theta)=\cos^{2}\theta+\sin^{2}\theta=1\ne 0$; su inversa es $Q_\theta\T$.
  \item No es invertible: la segunda fila es el doble de la primera, el rango es~1 y $\det=1\cdot 4-2\cdot 2=0$.
  \item No es invertible: $(1,0,-1,2)\ne 0$ está en $\Nuc(A)$, de modo que falla la condición~(iii).
  \item Es invertible: es diagonal y $\det=3\cdot 0{,}5\cdot(-1)=-1{,}5\ne 0$; su inversa es $\diag(1/3,\,2,\,-1)$.
\end{enumerate}
\resultado \textbf{invertibles: (b) y (e); singulares: (a), (c) y (d).}
\comprobacion $\diag(3,\,0{,}5,\,-1)\cdot\diag(1/3,\,2,\,-1)=\diag(1,1,1)$.
\end{solucion}

\begin{solucion}{eje:2-inversa}
\datos $A=\begin{pmatrix}1&2\\3&4\end{pmatrix}$.
\planteamiento \cref{prop:inversa2}.
\pasos
\begin{enumerate}[label=\arabic*.]
  \item $\det A=1\cdot 4-2\cdot 3=-2\ne 0$: $A$ es invertible.
  \item Matriz $\begin{pmatrix}d&-b\\-c&a\end{pmatrix}=\begin{pmatrix}4&-2\\-3&1\end{pmatrix}$.
  \item Dividiendo por $-2$: $A^{-1}=\begin{pmatrix}-2&1\\1{,}5&-0{,}5\end{pmatrix}$.
  \item Fila 1 de $AA^{-1}$: $(1\cdot(-2)+2\cdot 1{,}5;\ 1\cdot 1+2\cdot(-0{,}5))=(1,0)$.
  \item Fila 2 de $AA^{-1}$: $(3\cdot(-2)+4\cdot 1{,}5;\ 3\cdot 1+4\cdot(-0{,}5))=(0,1)$.
  \item $\det(A^{-1})=(-2)(-0{,}5)-1\cdot 1{,}5=1-1{,}5=-0{,}5$.
\end{enumerate}
\resultado \textbf{$A^{-1}=\begin{pmatrix}-2&1\\1{,}5&-0{,}5\end{pmatrix}$.}
\comprobacion $AA^{-1}=I$ por los pasos 4 y 5, y $\det(A^{-1})=-0{,}5=1/(-2)=1/\det A$.
\end{solucion}

\begin{solucion}{eje:2-mml}
\datos $A=\begin{pmatrix}1&3&5\\2&4&6\\0&2&4\end{pmatrix}$, $B=\begin{pmatrix}2&3&4\\3&4&5\\4&5&6\end{pmatrix}$.
\planteamiento \cref{prop:laplace} y regla de Sarrus para $A$; para $B$, calcular el determinante y, si es nulo, buscar un vector del núcleo.
\pasos
\begin{enumerate}[label=\arabic*.]
  \item Laplace por la primera fila de $A$: $1\cdot(4\cdot 4-6\cdot 2)-3\cdot(2\cdot 4-6\cdot 0)+5\cdot(2\cdot 2-4\cdot 0)=1\cdot 4-3\cdot 8+5\cdot 4=4-24+20=0$.
  \item Sarrus: diagonales descendentes $1\cdot 4\cdot 4+3\cdot 6\cdot 0+5\cdot 2\cdot 2=16+0+20=36$; ascendentes $5\cdot 4\cdot 0+1\cdot 6\cdot 2+3\cdot 2\cdot 4=0+12+24=36$. $\det A=36-36=0$.
  \item $\det B$ por la primera fila: $2\cdot(4\cdot 6-5\cdot 5)-3\cdot(3\cdot 6-5\cdot 4)+4\cdot(3\cdot 5-4\cdot 4)=2\cdot(-1)-3\cdot(-2)+4\cdot(-1)=-2+6-4=0$.
  \item $B$ no es invertible. Las filas están en progresión aritmética: $F_1-2F_2+F_3=(2-6+4,\,3-8+5,\,4-10+6)=0$. Lo mismo ocurre con las columnas, luego $B(1,-2,1)=0$.
\end{enumerate}
\resultado \textbf{(a) $\det A=0$ por ambos métodos. (b) $B$ no tiene inversa; $(1,-2,1)\in\Nuc(B)$.}
\comprobacion $B(1,-2,1)=(2-6+4,\,3-8+5,\,4-10+6)=(0,0,0)$. En $A$, la segunda columna es la media de la primera y la tercera: $\tfrac12\bigl((1,2,0)+(5,6,4)\bigr)=(3,4,2)$, lo que explica que $\det A=0$.
\end{solucion}

\begin{solucion}{eje:2-giros}
\datos $Q_\theta=\begin{pmatrix}\cos\theta&-\sin\theta\\ \sin\theta&\cos\theta\end{pmatrix}$.
\planteamiento multiplicar y usar las fórmulas del coseno y del seno de una suma; para (c), $Q_\alpha\T=Q_{-\alpha}$.
\pasos
\begin{enumerate}[label=\arabic*.]
  \item (a) Entrada $(1,1)$ de $Q_\alpha Q_\beta$: $\cos\alpha\cos\beta-\sin\alpha\sin\beta=\cos(\alpha+\beta)$. Entrada $(2,1)$: $\sin\alpha\cos\beta+\cos\alpha\sin\beta=\sin(\alpha+\beta)$.
  \item Entrada $(1,2)$: $-\cos\alpha\sin\beta-\sin\alpha\cos\beta=-\sin(\alpha+\beta)$. Entrada $(2,2)$: $-\sin\alpha\sin\beta+\cos\alpha\cos\beta=\cos(\alpha+\beta)$. Luego $Q_\alpha Q_\beta=Q_{\alpha+\beta}$.
  \item (b) $Q_{\pi/6}(2,0)=(2\cos 30^{\circ},\,2\sin 30^{\circ})=(\sqrt{3},\,1)\approx(1{,}732;\,1)$. Norma: $\sqrt{3+1}=2$.
  \item (c) $(Q_\alpha x)\T(Q_\beta y)=x\T Q_\alpha\T Q_\beta y=x\T Q_{-\alpha}Q_\beta y=x\T Q_{\beta-\alpha}y$, por el apartado~(a).
\end{enumerate}
\resultado \textbf{$Q_\alpha Q_\beta=Q_{\alpha+\beta}$; $Q_{\pi/6}(2,0)=(\sqrt{3},1)$, de norma~2; el producto escalar depende solo de $\beta-\alpha$.} Si $\alpha$ y $\beta$ son proporcionales a las posiciones de la consulta y de la clave, la puntuación depende solo de su diferencia.
\comprobacion con $\beta=-\alpha$, (a) da $Q_\alpha Q_{-\alpha}=Q_0=I$, coherente con $Q_\alpha^{-1}=Q_{-\alpha}$.
\end{solucion}

\begin{solucion}{eje:2-haar}
\datos $h_1=(1,1,1,1)$, $h_2=(1,1,-1,-1)$, $h_3=(1,-1,0,0)$, $h_4=(0,0,1,-1)$; $x=(4,2,5,5)$.
\planteamiento en una base ortogonal, $c_j=x\T h_j/h_j\T h_j$ (\cref{prop:coord-ortonormal}); el error de suprimir términos se calcula por Pitágoras.
\pasos
\begin{enumerate}[label=\arabic*.]
  \item (a) $h_1\T h_2=1+1-1-1=0$; $h_1\T h_3=1-1=0$; $h_1\T h_4=1-1=0$; $h_2\T h_3=1-1=0$; $h_2\T h_4=-1+1=0$; $h_3\T h_4=0$.
  \item Normas: $\norm{h_1}=\norm{h_2}=2$ y $\norm{h_3}=\norm{h_4}=\sqrt{2}$. Base ortonormal: $h_1/2$, $h_2/2$, $h_3/\sqrt{2}$, $h_4/\sqrt{2}$.
  \item (b) $x\T h_1=4+2+5+5=16$, $c_1=16/4=4$ (la media de la señal).
  \item $x\T h_2=4+2-5-5=-4$, $c_2=-4/4=-1$.
  \item $x\T h_3=4-2=2$, $c_3=2/2=1$; $x\T h_4=5-5=0$, $c_4=0$. Luego $x=4h_1-h_2+h_3$.
  \item (c) $\tilde x=4h_1-h_2=(4-1,\,4-1,\,4+1,\,4+1)=(3,3,5,5)$ y $x-\tilde x=(1,-1,0,0)$, con error cuadrático $1+1=2$.
  \item Justificación: $x-\tilde x=c_3h_3+c_4h_4$ y, por ortogonalidad, $\norm{x-\tilde x}^{2}=c_3^{2}\norm{h_3}^{2}+c_4^{2}\norm{h_4}^{2}$.
\end{enumerate}
\resultado \textbf{$x=4h_1-h_2+h_3$; la versión comprimida es $(3,3,5,5)$, con error cuadrático~2.}
\comprobacion $c_3^{2}\norm{h_3}^{2}+c_4^{2}\norm{h_4}^{2}=1\cdot 2+0\cdot 2=2$. Es la misma lógica del \cref{teo:eckart-young}: descartar términos cuesta exactamente su energía.
\end{solucion}

\begin{solucion}{eje:2-isometria}
\datos $Q,Q_1,Q_2\in\R^{n\times n}$ con $Q\T Q=Q_1\T Q_1=Q_2\T Q_2=I$.
\planteamiento usar $Q\T Q=I$ y las propiedades del determinante.
\pasos
\begin{enumerate}[label=\arabic*.]
  \item $(Qx)\T(Qy)=x\T Q\T Qy=x\T Iy=x\T y$.
  \item Tomando $y=x$: $\norm{Qx}^{2}=\norm{x}^{2}$ y, como ambas normas son no negativas, $\norm{Qx}=\norm{x}$.
  \item $1=\det(I)=\det(Q\T Q)=\det(Q\T)\det(Q)=\det(Q)^{2}$, luego $\det Q=\pm 1$.
  \item $(Q_1Q_2)\T(Q_1Q_2)=Q_2\T Q_1\T Q_1Q_2=Q_2\T Q_2=I$: el producto es ortogonal.
\end{enumerate}
\resultado \textbf{las matrices ortogonales conservan productos escalares y normas, tienen determinante $\pm 1$ y forman un conjunto cerrado para el producto.} \hfill$\blacksquare$
\comprobacion el recíproco de $\det Q=\pm1$ es falso: $\diag(2,\,0{,}5)$ tiene determinante~1 y no es ortogonal, porque $\norm{(2,0{,}5)e_1}=2\ne 1$.
\end{solucion}

% ---------------------------------------------------------------------
\section{Tema 3}

\begin{solucion}{eje:3-mml}
\datos $I_2$, $P=\diag(1,0)$, $J=\begin{pmatrix}1&1\\0&1\end{pmatrix}$, $N=\begin{pmatrix}0&1\\0&0\end{pmatrix}$.
\planteamiento invertible si y solo si $\det\ne 0$; diagonalizable si y solo si hay dos autovectores independientes (\cref{teo:diagonalizacion}).
\pasos
\begin{enumerate}[label=\arabic*.]
  \item $I_2$: ya es diagonal; $\det=1$. Diagonalizable e invertible.
  \item $P$: ya es diagonal; $\det=0$. Diagonalizable y no invertible.
  \item $J$: triangular con autovalor doble~1. $J-I=\begin{pmatrix}0&1\\0&0\end{pmatrix}$ da $v_2=0$, espacio propio $\gen\{e_1\}$ de dimensión~1. No diagonalizable; $\det=1$, invertible.
  \item $N$: autovalor doble~0; $Nv=0$ da $v_2=0$, espacio propio $\gen\{e_1\}$. No diagonalizable; $\det=0$, no invertible.
\end{enumerate}
\resultado \textbf{$I_2$: diagonalizable e invertible; $P$: diagonalizable, no invertible; $J$: invertible, no diagonalizable; $N$: ni lo uno ni lo otro.} Las cuatro combinaciones aparecen, luego ninguna propiedad implica la otra.
\comprobacion la invertibilidad depende de si $0$ es autovalor; la diagonalizabilidad, de cuántos autovectores independientes hay. Son preguntas distintas.
\end{solucion}

\begin{solucion}{eje:3-autovec}
\datos $B=\begin{pmatrix}4&1\\2&3\end{pmatrix}$.
\planteamiento resolver $\det(B-\lambda I)=0$ y, para cada raíz, $(B-\lambda I)v=0$.
\pasos
\begin{enumerate}[label=\arabic*.]
  \item $\det(B-\lambda I)=(4-\lambda)(3-\lambda)-2=\lambda^{2}-7\lambda+10$.
  \item $\lambda=(7\pm\sqrt{49-40})/2=(7\pm 3)/2$: $\lambda_1=5$, $\lambda_2=2$.
  \item $\lambda=5$: $B-5I=\begin{pmatrix}-1&1\\2&-2\end{pmatrix}$, ecuación $-v_1+v_2=0$, autovector $(1,1)$.
  \item $\lambda=2$: $B-2I=\begin{pmatrix}2&1\\2&1\end{pmatrix}$, ecuación $2v_1+v_2=0$, autovector $(1,-2)$.
  \item $(1,1)\T(1,-2)=1-2=-1\ne 0$: no son ortogonales.
\end{enumerate}
\resultado \textbf{$\lambda=5$ con $v=(1,1)$ y $\lambda=2$ con $v=(1,-2)$; los autovectores no son ortogonales porque $B$ no es simétrica ($b_{12}=1\ne b_{21}=2$)} y el \cref{teo:simetricas} no se aplica.
\comprobacion $B(1,1)=(5,5)$ y $B(1,-2)=(2,-4)$. Traza: $4+3=7=5+2$; determinante: $12-2=10=5\cdot 2$.
\end{solucion}

\begin{solucion}{eje:3-u2}
\datos las matrices $A$, $B$ y $C$ del enunciado.
\planteamiento $A$ es triangular; para $B$ se desarrolla por la fila central; $C$ tiene rango~1.
\pasos
\begin{enumerate}[label=\arabic*.]
  \item $A$ triangular superior: $\det(A-\lambda I)=(2-\lambda)(3-\lambda)(-1-\lambda)$, autovalores $2$, $3$ y $-1$.
  \item $B$: la fila~2 de $B-\lambda I$ es $(0,\,1-\lambda,\,0)$; desarrollando por ella, $\det(B-\lambda I)=(1-\lambda)\det\begin{pmatrix}-\lambda&2\\8&-\lambda\end{pmatrix}=(1-\lambda)(\lambda^{2}-16)$. Autovalores $1$, $4$ y $-4$.
  \item $C=3\,\mathbf{1}\mathbf{1}\T$ con $\mathbf{1}=(1,1,1)$ tiene rango~1, luego $\dim\Nuc(C)=2$ y el autovalor~0 tiene dos autovectores independientes. El tercero es $C\mathbf{1}=(9,9,9)=9\cdot\mathbf{1}$. Autovalores $9$, $0$ y $0$.
  \item Trazas: $A$: $2+3-1=4$; $B$: $0+1+0=1=1+4-4$; $C$: $3+3+3=9=9+0+0$.
  \item Determinantes: $A$: $2\cdot 3\cdot(-1)=-6$; $B$: desarrollando por la fila~2, $1\cdot(0\cdot 0-2\cdot 8)=-16=1\cdot 4\cdot(-4)$; $C$: $0=9\cdot 0\cdot 0$.
\end{enumerate}
\resultado \textbf{$A$: $2$, $3$, $-1$. $B$: $1$, $4$, $-4$. $C$: $9$, $0$, $0$.}
\comprobacion los pasos 4 y 5 verifican la \cref{prop:traza-det} en las tres matrices.
\end{solucion}

\begin{solucion}{eje:3-u3}
\datos autovalores de $B$: $0$, $2$ y $3$, distintos.
\planteamiento autovalores distintos implican diagonalizable (\cref{prop:indep-autovectores}); si $Bv=\lambda v$, entonces $p(B)v=p(\lambda)v$ para todo polinomio $p$.
\pasos
\begin{enumerate}[label=\arabic*.]
  \item (a) $B=X\Lambda X^{-1}$ con $\Lambda=\diag(0,2,3)$. Multiplicar por matrices invertibles no cambia el rango, luego $\rango(B)=\rango(\Lambda)=2$.
  \item (b) $\det B=0\cdot 2\cdot 3=0$ y $\det(B\T B)=\det(B\T)\det(B)=0$.
  \item (c) Si $Bv=\lambda v$: $(B^{2}+I)v=(\lambda^{2}+1)v$, luego $(B^{2}+I)^{-1}v=v/(\lambda^{2}+1)$.
  \item Con $\lambda=0,2,3$: $1/1=1$, $1/5=0{,}2$ y $1/10=0{,}1$.
  \item (d) No. $B_1=\diag(0,2,3)$ da $B_1\T B_1=\diag(0,4,9)$. La matriz triangular $B_2=\begin{pmatrix}0&1&0\\0&2&0\\0&0&3\end{pmatrix}$ tiene los mismos autovalores y $B_2\T B_2=\diag(0,5,9)$, con autovalores $0$, $5$ y $9$.
\end{enumerate}
\resultado \textbf{(a) rango~2; (b) 0; (c) $1$, $1/5$ y $1/10$; (d) no pueden deducirse.}
\comprobacion $B^{2}+I$ es invertible porque $\lambda^{2}+1\ge 1$. En (d), $B_2$ no es simétrica; solo para matrices simétricas se cumple $\sigma_i=\abs{\lambda_i}$ (\cref{cor:svd-props}).
\end{solucion}

\begin{solucion}{eje:3-reales}
\datos $S\in\R^{n\times n}$ con $S\T=S$; $Sv=\lambda v$ con $v\in\mathbb{C}^{n}$, $v\ne 0$.
\planteamiento calcular el número $\bar v\T Sv$ de dos formas.
\pasos
\begin{enumerate}[label=\arabic*.]
  \item Conjugando $Sv=\lambda v$ y usando que $S$ es real: $S\bar v=\bar\lambda\bar v$.
  \item Primera forma: $\bar v\T(Sv)=\lambda\,\bar v\T v$.
  \item Segunda forma: $\bar v\T Sv=(S\T\bar v)\T v=(S\bar v)\T v=\bar\lambda\,\bar v\T v$, por la simetría y el paso~1.
  \item Restando: $(\lambda-\bar\lambda)\,\bar v\T v=0$.
  \item $\bar v\T v=\sum_i\abs{v_i}^{2}>0$ porque $v\ne 0$, luego $\lambda=\bar\lambda$ y $\lambda\in\R$.
\end{enumerate}
\resultado \textbf{todo autovalor de una matriz real simétrica es real.} Como $S-\lambda I$ es entonces una matriz real singular, también pueden tomarse autovectores reales. \hfill$\blacksquare$
\comprobacion sin simetría el resultado falla: el giro $\begin{pmatrix}0&-1\\1&0\end{pmatrix}$ tiene autovalores $\pm i$ (\cref{tab:autovalores}).
\end{solucion}

\begin{solucion}{eje:3-potencia}
\datos $Sq_i=\lambda_iq_i$ con $q_i$ ortonormales; $\abs{\lambda_1}>\abs{\lambda_2}\ge\dots\ge\abs{\lambda_n}$; $x_0=\sum_ic_iq_i$ con $c_1\ne 0$.
\planteamiento en la base de autovectores, $S^{k}$ multiplica cada coordenada por $\lambda_i^{k}$.
\pasos
\begin{enumerate}[label=\arabic*.]
  \item Por inducción, $S^{k}q_i=\lambda_i^{k}q_i$; por linealidad, $S^{k}x_0=\sum_ic_i\lambda_i^{k}q_i$.
  \item Sacando factor común $c_1\lambda_1^{k}\ne 0$: $S^{k}x_0=c_1\lambda_1^{k}(q_1+e_k)$, con $e_k=\sum_{i\ge 2}(c_i/c_1)(\lambda_i/\lambda_1)^{k}q_i$.
  \item Por ortonormalidad, $\norm{e_k}^{2}=\sum_{i\ge 2}(c_i/c_1)^{2}(\lambda_i/\lambda_1)^{2k}\le(\lambda_2/\lambda_1)^{2k}\sum_{i\ge 2}(c_i/c_1)^{2}$.
  \item Luego $\norm{e_k}\le K\abs{\lambda_2/\lambda_1}^{k}\to 0$, con $K=\bigl(\sum_{i\ge 2}c_i^{2}\bigr)^{1/2}/\abs{c_1}$.
  \item Normalizando: $S^{k}x_0/\norm{S^{k}x_0}=s_k\,(q_1+e_k)/\norm{q_1+e_k}$, con $s_k=\operatorname{signo}(c_1\lambda_1^{k})=\pm 1$. Como $e_k\perp q_1$, $\norm{q_1+e_k}=(1+\norm{e_k}^{2})^{1/2}\to 1$.
  \item La componente de $x_k$ ortogonal a $q_1$ es $s_k\,e_k/\norm{q_1+e_k}$, de norma a lo sumo $\norm{e_k}\le K\abs{\lambda_2/\lambda_1}^{k}$.
\end{enumerate}
\resultado \textbf{$S^{k}x_0/\norm{S^{k}x_0}$ se aproxima a $\pm q_1$ (el signo alterna si $\lambda_1<0$), y su distancia a la recta $\gen\{q_1\}$ es $O(\abs{\lambda_2/\lambda_1}^{k})$.} \hfill$\blacksquare$
\comprobacion en el \cref{ej:potencia}, $x_0=(q_1+q_2)/\sqrt{2}$, de modo que $c_1=c_2$ y $K=1$; tras 5 pasos, la cota es $3^{-5}\approx 0{,}004115$, y la componente de $S^{5}x_0/\norm{S^{5}x_0}$ sobre $q_2$ vale $(0{,}7100-0{,}7042)/\sqrt{2}\approx 0{,}0041$.
\end{solucion}

% ---------------------------------------------------------------------
\section{Tema 4}

\begin{solucion}{eje:4-vf}
\planteamiento contrastar cada afirmación con el \cref{teo:svd} y el \cref{cor:svd-props}.
\pasos
\begin{enumerate}[label=(\alph*)]
  \item \textbf{Falsa.} El \cref{teo:svd} vale para cualquier $A\in\R^{m\times n}$; el \cref{ej:svd23} trata una matriz de $2\times 3$.
  \item \textbf{Falsa.} $\sigma_i=\sqrt{\lambda_i(A\T A)}\ge 0$ por definición. Un signo negativo puede trasladarse siempre a $u_i$.
  \item \textbf{Falsa en general.} Son los valores absolutos de los autovalores: $\diag(1,-2)$ tiene autovalores $1$ y $-2$ y valores singulares $2$ y $1$. Es verdadera si además $A$ es semidefinida positiva.
  \item \textbf{Verdadera.} Si $A=U\Sigma V\T$, entonces $A\T=V\Sigma\T U\T$ es una SVD de $A\T$ con los mismos $\sigma_i$.
\end{enumerate}
\resultado \textbf{(a) falsa; (b) falsa; (c) falsa salvo que $A$ sea semidefinida positiva; (d) verdadera.}
\end{solucion}

\begin{solucion}{eje:4-u1}
\datos $A=\begin{pmatrix}1&1\\1&0\end{pmatrix}$.
\planteamiento $\sigma_i=\sqrt{\lambda_i(A\T A)}$.
\pasos
\begin{enumerate}[label=\arabic*.]
  \item (a) $A$ es simétrica, luego $A\T A=AA\T=A^{2}$. Fila 1: $(1+1,\,1+0)=(2,1)$; fila 2: $(1+0,\,1+0)=(1,1)$. $A\T A=\begin{pmatrix}2&1\\1&1\end{pmatrix}$.
  \item (b) $\det(A\T A-\lambda I)=(2-\lambda)(1-\lambda)-1=\lambda^{2}-3\lambda+1$, con raíces $\lambda=(3\pm\sqrt{5})/2$, es decir, $2{,}618$ y $0{,}3820$.
  \item $\sigma_1=\sqrt{2{,}618}=1{,}618$ y $\sigma_2=\sqrt{0{,}3820}=0{,}6180$.
  \item (c) $\bigl((1+\sqrt{5})/2\bigr)^{2}=(6+2\sqrt{5})/4=(3+\sqrt{5})/2$ y $\bigl((\sqrt{5}-1)/2\bigr)^{2}=(6-2\sqrt{5})/4=(3-\sqrt{5})/2$.
  \item Los autovalores de $A$ son las raíces de $\lambda^{2}-\lambda-1=0$, $(1\pm\sqrt{5})/2\approx 1{,}618$ y $-0{,}6180$. Como $A$ es simétrica, el \cref{cor:svd-props}(iv) da $\sigma_i=\abs{\lambda_i}$.
\end{enumerate}
\resultado \textbf{$\sigma_1=(1+\sqrt{5})/2\approx 1{,}618$ y $\sigma_2=(\sqrt{5}-1)/2\approx 0{,}6180$}, los valores absolutos de los autovalores de $A$, porque $A$ es simétrica.
\comprobacion $\sigma_1\sigma_2=(5-1)/4=1=\abs{\det A}$ y $\sigma_1^{2}+\sigma_2^{2}=3=\norm{A}_F^{2}=1+1+1+0$.
\end{solucion}

\begin{solucion}{eje:4-u5}
\datos columnas ortogonales $w_j$ con $\norm{w_j}=s_j>0$; $A=\begin{pmatrix}2&1\\1&0\\2&-1\end{pmatrix}$.
\planteamiento $(A\T A)_{ij}=w_i\T w_j$; comparar con $A\T A=V\Sigma\T\Sigma V\T$.
\pasos
\begin{enumerate}[label=\arabic*.]
  \item $(A\T A)_{ij}=w_i\T w_j$ vale $0$ si $i\ne j$ y $s_i^{2}$ si $i=j$: $A\T A=\diag(s_1^{2},\dots,s_n^{2})$.
  \item $A\T A$ ya es diagonal: $V=I$ (con las columnas ordenadas de mayor a menor norma), $\Sigma=\diag(s_1,\dots,s_n)$ y, de $AV=U\Sigma$, $u_j=w_j/s_j$.
  \item Caso numérico: $w_1=(2,1,2)$ y $w_2=(1,0,-1)$, con $w_1\T w_2=2+0-2=0$.
  \item $s_1=\sqrt{4+1+4}=3$ y $s_2=\sqrt{1+0+1}=\sqrt{2}\approx 1{,}414$; $A\T A=\diag(9,2)$.
  \item $u_1=(2,1,2)/3$ y $u_2=(1,0,-1)/\sqrt{2}$; para la SVD completa se añade $u_3$ unitario y ortogonal a ambos, $u_3=(1,-4,1)/\sqrt{18}$.
\end{enumerate}
\resultado \textbf{$A\T A=\diag(s_j^{2})$, $V=I$, $\Sigma=\diag(s_j)$ y $U$ tiene por columnas las $w_j$ normalizadas. En el ejemplo, $\sigma=(3,\sqrt{2})$.}
\comprobacion $[\,3u_1\mid\sqrt{2}\,u_2\,]=[(2,1,2)\mid(1,0,-1)]=A$ y $\norm{A}_F^{2}=4+1+1+0+4+1=11=9+2$. Además, $u_3\T u_1=(2-4+2)/(3\sqrt{18})=0$ y $u_3\T u_2=(1-1)/6=0$.
\end{solucion}

\begin{solucion}{eje:4-mml}
\datos $A=\begin{pmatrix}2&2\\-1&1\end{pmatrix}$.
\planteamiento autovectores de $A\T A$ y $u_i=Av_i/\sigma_i$.
\pasos
\begin{enumerate}[label=\arabic*.]
  \item $A\T A=\begin{pmatrix}2&-1\\2&1\end{pmatrix}\begin{pmatrix}2&2\\-1&1\end{pmatrix}=\begin{pmatrix}4+1&4-1\\4-1&4+1\end{pmatrix}=\begin{pmatrix}5&3\\3&5\end{pmatrix}$.
  \item $(5-\lambda)^{2}=9$, luego $\lambda=8$ y $\lambda=2$: $\sigma_1=\sqrt{8}=2\sqrt{2}\approx 2{,}828$ y $\sigma_2=\sqrt{2}\approx 1{,}414$.
  \item Autovectores: $v_1=(1,1)/\sqrt{2}$ para $\lambda=8$ y $v_2=(-1,1)/\sqrt{2}$ para $\lambda=2$.
  \item $u_1=Av_1/\sigma_1=(4,0)/(\sqrt{2}\cdot 2\sqrt{2})=(1,0)$.
  \item $u_2=Av_2/\sigma_2=(-2+2,\,1+1)/(\sqrt{2}\cdot\sqrt{2})=(0,2)/2=(0,1)$.
\end{enumerate}
\resultado \textbf{$U=I_2$, $\Sigma=\diag(2\sqrt{2},\sqrt{2})$, $V=\frac{1}{\sqrt{2}}\begin{pmatrix}1&-1\\1&1\end{pmatrix}$.}
\comprobacion $U\Sigma V\T=\begin{pmatrix}2\sqrt{2}&0\\0&\sqrt{2}\end{pmatrix}\frac{1}{\sqrt{2}}\begin{pmatrix}1&1\\-1&1\end{pmatrix}=\begin{pmatrix}2&2\\-1&1\end{pmatrix}=A$. También $\sigma_1\sigma_2=4=\abs{\det A}=\abs{2+2}$.
\end{solucion}

\begin{solucion}{eje:4-frobenius}
\datos $A\in\R^{m\times n}$ con SVD $A=U\Sigma V\T$.
\planteamiento expresar $\norm{A}_F^{2}$ como una traza y usar $\traza(XY)=\traza(YX)$.
\pasos
\begin{enumerate}[label=\arabic*.]
  \item $(A\T A)_{jj}=\sum_ia_{ij}^{2}$; sumando en $j$, $\traza(A\T A)=\sum_{i,j}a_{ij}^{2}=\norm{A}_F^{2}$.
  \item $A\T A=V\Sigma\T U\T U\Sigma V\T=V(\Sigma\T\Sigma)V\T$, porque $U\T U=I$.
  \item $\traza(V(\Sigma\T\Sigma)V\T)=\traza((\Sigma\T\Sigma)V\T V)=\traza(\Sigma\T\Sigma)$, porque $V\T V=I$.
  \item $\Sigma\T\Sigma=\diag(\sigma_1^{2},\dots,\sigma_r^{2},0,\dots,0)$, de traza $\sum_i\sigma_i^{2}$.
\end{enumerate}
\resultado \textbf{$\norm{A}_F^{2}=\traza(A\T A)=\sum_i\sigma_i^{2}$.} \hfill$\blacksquare$
\comprobacion para el \cref{ej:svd23}: $9+4+4+4+9+4=34=25+9$.
\end{solucion}

\begin{solucion}{eje:4-autovalores}
\datos $A\in\R^{m\times n}$.
\planteamiento transportar autovectores de $A\T A$ a $AA\T$ mediante $v\mapsto Av$ y viceversa.
\pasos
\begin{enumerate}[label=\arabic*.]
  \item Sea $A\T Av=\lambda v$ con $v\ne 0$ y $\lambda\ne 0$. Entonces $Av\ne 0$, pues en otro caso $\lambda v=A\T(Av)=0$.
  \item $AA\T(Av)=A(A\T Av)=\lambda\,Av$: $\lambda$ es autovalor de $AA\T$ con autovector $Av$.
  \item Multiplicidades: la aplicación $v\mapsto Av$ es inyectiva sobre el espacio propio $E_\lambda$ de $A\T A$, porque si $Av=0$ con $v\in E_\lambda$, entonces $\lambda v=0$ y $v=0$. Lleva $E_\lambda$ al espacio propio de $\lambda$ en $AA\T$, luego $\dim E_\lambda(A\T A)\le\dim E_\lambda(AA\T)$.
  \item Intercambiando los papeles de $A$ y $A\T$ se obtiene la desigualdad contraria.
  \item Ambas matrices son simétricas y, por tanto, diagonalizables (\cref{teo:espectral}): la multiplicidad de cada autovalor es la dimensión de su espacio propio.
\end{enumerate}
\resultado \textbf{$A\T A$ y $AA\T$ tienen los mismos autovalores no nulos con las mismas multiplicidades}; difieren solo en el número de ceros, $n-r$ frente a $m-r$. \hfill$\blacksquare$
\comprobacion en el \cref{ej:svd23}, $AA\T$ tiene autovalores $25$ y $9$, y $A\T A$, de orden~3, tiene $25$, $9$ y $0$.
\end{solucion}

\begin{solucion}{eje:4-espectral}
\datos $A\in\R^{m\times n}$ con SVD $A=\sum_i\sigma_iu_iv_i\T$.
\planteamiento escribir $x$ en la base ortonormal $\{v_i\}$ de $\R^{n}$.
\pasos
\begin{enumerate}[label=\arabic*.]
  \item $x=\sum_{i=1}^{n}c_iv_i$ con $\norm{x}^{2}=\sum_ic_i^{2}$ (\cref{prop:coord-ortonormal}).
  \item $Ax=\sum_{i\le r}c_i\sigma_iu_i$ y, por ortonormalidad de los $u_i$, $\norm{Ax}^{2}=\sum_{i\le r}c_i^{2}\sigma_i^{2}$.
  \item $\sum_{i\le r}c_i^{2}\sigma_i^{2}\le\sigma_1^{2}\sum_ic_i^{2}=\sigma_1^{2}\norm{x}^{2}$, luego $\norm{Ax}/\norm{x}\le\sigma_1$.
  \item Con $x=v_1$: $Av_1=\sigma_1u_1$ y $\norm{Av_1}/\norm{v_1}=\sigma_1$.
\end{enumerate}
\resultado \textbf{$\norm{A}_2=\max_{x\ne 0}\norm{Ax}/\norm{x}=\sigma_1$, alcanzado en $v_1$.} \hfill$\blacksquare$
\comprobacion en el \cref{eje:4-mml}, $\norm{Av_1}=\norm{(4,0)/\sqrt{2}}=2\sqrt{2}=\sigma_1$, y $e_2$ da $\norm{Ae_2}=\norm{(2,1)}=\sqrt{5}\approx 2{,}236<2{,}828$.
\end{solucion}

% ---------------------------------------------------------------------
\section{Tema 5}

\begin{solucion}{eje:5-centrar}
\datos $n$ puntos $x_i$ con media $\mu$ y covarianza $C=\frac1n\sum_i(x_i-\mu)(x_i-\mu)\T$.
\planteamiento PCA busca la dirección de mayor dispersión alrededor de la media; se compara $C$ con la matriz sin centrar $\frac1n\sum_ix_ix_i\T$.
\pasos
\begin{enumerate}[label=\arabic*.]
  \item Escribiendo $x_i=(x_i-\mu)+\mu$: $\frac1n\sum_ix_ix_i\T=C+\frac1n\sum_i\bigl[(x_i-\mu)\mu\T+\mu(x_i-\mu)\T\bigr]+\mu\mu\T$.
  \item Los términos cruzados se anulan porque $\sum_i(x_i-\mu)=0$: $\frac1n\sum_ix_ix_i\T=C+\mu\mu\T$.
  \item Si los datos están lejos del origen, $\norm{\mu}^{2}$ domina sobre los autovalores de $C$, y $\mu\mu\T$ es de rango~1 con autovector $\mu/\norm{\mu}$.
  \item La primera «componente» apunta entonces hacia la media, es decir, describe la posición de la nube y no la dirección en la que se estira.
\end{enumerate}
\resultado \textbf{hay que centrar porque PCA debe medir dispersión, no posición; sin centrar, la primera componente se aproxima a $\mu/\norm{\mu}$.}
\comprobacion con los puntos $(101,99)$ y $(99,101)$, la nube se estira en la dirección $(1,-1)$: centrada, $C=\begin{pmatrix}1&-1\\-1&1\end{pmatrix}$, con autovector principal $(1,-1)/\sqrt{2}$. Sin centrar, $\frac12\sum_ix_ix_i\T=\begin{pmatrix}10\,001&9999\\9999&10\,001\end{pmatrix}$, cuyo autovector principal es $(1,1)/\sqrt{2}$, con autovalor $20\,000$.
\end{solucion}

\begin{solucion}{eje:5-varexp}
\datos varianzas explicadas $0{,}12$ y $0{,}07$; dimensión $768$.
\planteamiento las varianzas explicadas se suman; una proyección ortogonal no alarga distancias.
\pasos
\begin{enumerate}[label=\arabic*.]
  \item Fracción mostrada: $0{,}12+0{,}07=0{,}19$, el $19\,\%$. El $81\,\%$ restante queda oculto.
  \item Sea $P$ la proyección ortogonal sobre el plano de las dos componentes. Por Pitágoras, $\norm{x-y}^{2}=\norm{P(x-y)}^{2}+\norm{(I-P)(x-y)}^{2}\ge\norm{P(x-y)}^{2}$.
  \item Lejos en el gráfico implica lejos en $\R^{768}$: la distancia real es al menos la del gráfico.
  \item Cerca en el gráfico no implica nada: los puntos pueden diferir mucho en las direcciones descartadas.
\end{enumerate}
\resultado \textbf{el gráfico muestra el $19\,\%$ de la varianza. Las separaciones son fiables; las cercanías, no.}
\comprobacion $(0,0,5)$ y $(0,0,-5)$ se proyectan sobre el mismo punto del plano $xy$ y están a distancia~10.
\end{solucion}

\begin{solucion}{eje:5-svd}
\datos $A=\begin{pmatrix}3&0\\4&5\end{pmatrix}$.
\planteamiento $\sigma_i=\sqrt{\lambda_i(A\T A)}$; por el \cref{teo:eckart-young}, $\norm{A-A_1}_F=\sigma_2$; energía $\sigma_1^{2}/(\sigma_1^{2}+\sigma_2^{2})$.
\pasos
\begin{enumerate}[label=\arabic*.]
  \item $A\T A$: entrada $(1,1)=9+16=25$, $(1,2)=0+20=20$, $(2,2)=0+25=25$. $A\T A=\begin{pmatrix}25&20\\20&25\end{pmatrix}$.
  \item $(25-\lambda)^{2}=400$, luego $\lambda=45$ y $\lambda=5$.
  \item (a) $\sigma_1=\sqrt{45}=3\sqrt{5}\approx 6{,}708$ y $\sigma_2=\sqrt{5}\approx 2{,}236$.
  \item (b) $\norm{A-A_1}_F=\sigma_2=\sqrt{5}\approx 2{,}236$.
  \item (c) Energía: $45/(45+5)=0{,}9$.
  \item (d) $v_1=(1,1)/\sqrt{2}$; $Av_1=(3,9)/\sqrt{2}$; $u_1=Av_1/\sigma_1=(3,9)/(\sqrt{2}\cdot 3\sqrt{5})=(1,3)/\sqrt{10}$.
  \item $A_1=\sigma_1u_1v_1\T=3\sqrt{5}\cdot\frac{1}{\sqrt{20}}\begin{pmatrix}1&1\\3&3\end{pmatrix}=\frac32\begin{pmatrix}1&1\\3&3\end{pmatrix}$.
\end{enumerate}
\resultado \textbf{(a) $\sigma_1=3\sqrt{5}\approx 6{,}708$, $\sigma_2=\sqrt{5}\approx 2{,}236$; (b) $\approx 2{,}236$; (c) $90\,\%$; (d) $A_1=\begin{pmatrix}1{,}5&1{,}5\\4{,}5&4{,}5\end{pmatrix}$.}
\comprobacion $A-A_1=\begin{pmatrix}1{,}5&-1{,}5\\-0{,}5&0{,}5\end{pmatrix}$, con $\norm{A-A_1}_F^{2}=2{,}25+2{,}25+0{,}25+0{,}25=5$. Además, $\norm{A}_F^{2}=9+16+25=50=45+5$ y $\sigma_1\sigma_2=15=\abs{\det A}$.
\end{solucion}

\begin{solucion}{eje:5-mml}
\datos SVD del \cref{eje:4-mml}: $\sigma_1=2\sqrt{2}$, $\sigma_2=\sqrt{2}$, $u_1=(1,0)$, $v_1=(1,1)/\sqrt{2}$.
\planteamiento $A_1=\sigma_1u_1v_1\T$ y \cref{teo:eckart-young}.
\pasos
\begin{enumerate}[label=\arabic*.]
  \item $A_1=2\sqrt{2}\begin{pmatrix}1\\0\end{pmatrix}\frac{1}{\sqrt{2}}\begin{pmatrix}1&1\end{pmatrix}=\begin{pmatrix}2&2\\0&0\end{pmatrix}$.
  \item $A-A_1=\begin{pmatrix}0&0\\-1&1\end{pmatrix}$, con $\norm{A-A_1}_F=\sqrt{2}$.
  \item $A-A_1=\sigma_2u_2v_2\T$ tiene un único valor singular, luego $\norm{A-A_1}_2=\sigma_2=\sqrt{2}$.
  \item Energía retenida: $\sigma_1^{2}/(\sigma_1^{2}+\sigma_2^{2})=8/10=0{,}8$.
\end{enumerate}
\resultado \textbf{$A_1=\begin{pmatrix}2&2\\0&0\end{pmatrix}$; error $\sqrt{2}\approx 1{,}414$ en ambas normas; energía retenida $80\,\%$.}
\comprobacion $\norm{A}_F^{2}=4+4+1+1=10=8+2$ y $\norm{A_1}_F^{2}=8=\sigma_1^{2}$.
\end{solucion}

\begin{solucion}{eje:5-pca}
\datos puntos $(4,2)$, $(2,4)$, $(-2,-4)$, $(-4,-2)$; $n=4$.
\planteamiento $C=X\T X/n$; autovalores y autovectores de $C$; proyecciones $z_i=x_i\T q_1$.
\pasos
\begin{enumerate}[label=\arabic*.]
  \item Media: $\bigl((4+2-2-4)/4,\,(2+4-4-2)/4\bigr)=(0,0)$; los datos ya están centrados.
  \item $\sum x_i^{2}=16+4+4+16=40$, $\sum y_i^{2}=4+16+16+4=40$ y $\sum x_iy_i=8+8+8+8=32$.
  \item $C=\frac14\begin{pmatrix}40&32\\32&40\end{pmatrix}=\begin{pmatrix}10&8\\8&10\end{pmatrix}$.
  \item $(10-\lambda)^{2}=64$: $\lambda_1=18$ y $\lambda_2=2$.
  \item $\lambda_1=18$: $C-18I=\begin{pmatrix}-8&8\\8&-8\end{pmatrix}$, luego $q_1=(1,1)/\sqrt{2}$. Varianza explicada: $18/20=0{,}9$.
  \item Proyecciones $z=(x+y)/\sqrt{2}$: $6/\sqrt{2}=3\sqrt{2}\approx 4{,}243$ para los dos primeros puntos y $-3\sqrt{2}\approx -4{,}243$ para los dos últimos.
\end{enumerate}
\resultado \textbf{$C=\begin{pmatrix}10&8\\8&10\end{pmatrix}$; $\lambda=18$ y $2$; $q_1=(1,1)/\sqrt{2}$ explica el $90\,\%$; proyecciones $(4{,}243;\,4{,}243;\,-4{,}243;\,-4{,}243)$.}
\comprobacion varianza de las proyecciones: $4\cdot 18/4=18=\lambda_1$. Por el \cref{teo:svd-pca}, los valores singulares de $X$ deben ser $\sqrt{4\cdot 18}=\sqrt{72}\approx 8{,}485$ y $\sqrt{4\cdot 2}=\sqrt{8}\approx 2{,}828$, como confirma el cálculo numérico.
\end{solucion}

\begin{solucion}{eje:5-coste}
\datos (a) $m\times n=1000\times 800$, $k=50$; (b) $d_{\text{in}}=d_{\text{out}}=4096$, $r=16$.
\planteamiento $A_k$ cuesta $k(m+n+1)$ números; LoRA entrena $r(d_{\text{in}}+d_{\text{out}})$ parámetros.
\pasos
\begin{enumerate}[label=\arabic*.]
  \item (a) $50\cdot(1000+800+1)=50\cdot 1801=90\,050$ números.
  \item Original: $1000\cdot 800=800\,000$. Fracción: $90\,050/800\,000=0{,}1126\approx 11{,}3\,\%$.
  \item Umbral: $1801\,k<800\,000\iff k<444{,}2$. Con $k=444$: $799\,644<800\,000$; con $k=445$: $801\,445>800\,000$.
  \item (b) $16\cdot(4096+4096)=16\cdot 8192=131\,072$ parámetros.
  \item $4096^{2}=16\,777\,216$; porcentaje: $131\,072/16\,777\,216=0{,}0078125=0{,}78\,\%$.
  \item Por la \cref{prop:rango-producto}, $\rango(\Delta W)=\rango(BA)\le 16$.
\end{enumerate}
\resultado \textbf{(a) $90\,050$ números ($11{,}3\,\%$); deja de compensar a partir de $k=445$. (b) $131\,072$ parámetros ($0{,}78\,\%$); $\rango(\Delta W)\le 16$.}
\comprobacion (b) duplica los valores de la \cref{obs:lora-rango} ($65\,536$ y $0{,}39\,\%$ con $r=8$), como corresponde a duplicar $r$. En (a), el rango máximo de la imagen es $800>445$: la compresión deja de compensar antes de ser exacta.
\end{solucion}

\begin{solucion}{eje:5-lora}
\datos $\Delta W^{\star}=\sum_i\sigma_iu_iv_i\T$; matrices $B\in\R^{d\times r}$, $A\in\R^{r\times n}$.
\planteamiento el conjunto de productos $BA$ coincide con el de matrices de rango a lo sumo $r$; después se aplica el \cref{teo:eckart-young}.
\pasos
\begin{enumerate}[label=\arabic*.]
  \item Por la \cref{prop:rango-producto}, todo producto $BA$ tiene rango $\le r$.
  \item Por el \cref{teo:eckart-young}, para toda matriz de rango $\le r$, $\norm{\Delta W^{\star}-BA}_F\ge(\sum_{i>r}\sigma_i^{2})^{1/2}$.
  \item La cota se alcanza: $B=U_r\Sigma_r=[\sigma_1u_1\mid\cdots\mid\sigma_ru_r]$ y $A=V_r\T$ dan $BA=\sum_{i\le r}\sigma_iu_iv_i\T=\Delta W^{\star}_r$.
  \item Caso numérico: $\diag(3,2,1)$ tiene valores singulares $3$, $2$ y $1$ con $u_i=v_i=e_i$. Con $r=1$, $B=3e_1$, $A=e_1\T$, $BA=\diag(3,0,0)$.
  \item Error: $\norm{\diag(0,2,1)}_F=\sqrt{4+1}=\sqrt{5}\approx 2{,}236$.
\end{enumerate}
\resultado \textbf{el mínimo es $(\sum_{i>r}\sigma_i^{2})^{1/2}$, alcanzado por $B=U_r\Sigma_r$, $A=V_r\T$; es nulo solo si $\rango(\Delta W^{\star})\le r$. En el ejemplo, $\sqrt{5}\approx 2{,}236$.} \hfill$\blacksquare$
\comprobacion la factorización no es única: $B=U_r$ y $A=\Sigma_rV_r\T$, o $BG$ y $G^{-1}A$ con cualquier $G\in\R^{r\times r}$ invertible, dan el mismo producto. En el ejemplo, $B=e_1$ y $A=3e_1\T$ también alcanzan el mínimo.
\end{solucion}

% ---------------------------------------------------------------------
\section{Tema 6}

\begin{solucion}{eje:6-formas}
\planteamiento \texttt{@} contrae el último eje del primer operando con el penúltimo del segundo; el \emph{broadcasting} alinea por la derecha (\cref{def:broadcasting}).
\pasos
\begin{enumerate}[label=\arabic*.]
  \item (a) $512=512$: el resultado es $(16,128,2048)$.
  \item (b) Se alinea $(16,128,512)$ con $(\cdot,128,1)$. Último eje: $512$ frente a $1$, se repite. Penúltimo: $128=128$. Primero: falta, se trata como~1 y se repite 16 veces. Resultado $(16,128,512)$.
  \item (c) Se alinea $(16,128,512)$ con $(\cdot,16,128)$. Último eje: $512$ frente a $128$, ni iguales ni~1: error.
  \item En (b), cada posición $t$ recibe un único número, el mismo para sus 512 rasgos. Un sesgo por rasgo debe tener forma $(512,)$.
\end{enumerate}
\resultado \textbf{(a) $(16,128,2048)$; (b) $(16,128,512)$, pero suma un valor por posición, no por rasgo; (c) error.}
\comprobacion en NumPy, (a) y (b) dan esas formas y (c) produce el error «operands could not be broadcast together».
\end{solucion}

\begin{solucion}{eje:6-d2l}
\datos \texttt{X} de forma $(2,3,4)$; \texttt{A} de forma $(3,3)$ o $(2,3)$.
\planteamiento una suma a lo largo de un eje lo elimina (o lo deja con tamaño~1 con \texttt{keepdims}); después se aplica la \cref{def:broadcasting}.
\pasos
\begin{enumerate}[label=\arabic*.]
  \item (a) Eje 0: $(3,4)$; eje 1: $(2,4)$; eje 2: $(2,3)$. Con \texttt{keepdims=True}: $(1,3,4)$, $(2,1,4)$ y $(2,3,1)$.
  \item (b) \texttt{A.sum(axis=1)} tiene forma $(3,)$ y contiene las sumas de las filas, $s_i=\sum_ja_{ij}$.
  \item Al dividir, $(3,3)/(3,)$ alinea el vector con el \emph{último} eje: el elemento $(i,j)$ se divide por $s_j$, la suma de la fila $j$, no por $s_i$.
  \item Corrección: \texttt{A / A.sum(axis=1, keepdims=True)}, de forma $(3,3)/(3,1)$, divide el elemento $(i,j)$ por $s_i$.
  \item Si \texttt{A} es de $(2,3)$, \texttt{A.sum(axis=1)} es de $(2,)$ y $(2,3)/(2,)$ compara $3$ con $2$ en el último eje: error.
\end{enumerate}
\resultado \textbf{(a) $(3,4)$, $(2,4)$, $(2,3)$. (b) divide cada columna $j$ por la suma de la fila $j$; la versión correcta conserva el eje con tamaño~1; con forma $(2,3)$ se produce un error.}
\comprobacion con $A=\begin{pmatrix}1&1&0\\0&2&2\\1&1&1\end{pmatrix}$, las sumas de filas son $(2,4,3)$; la versión correcta da la primera fila $(0{,}5;\,0{,}5;\,0)$, que suma~1, y la incorrecta da $(1/2,\,1/4,\,0)$, que suma $0{,}75$.
\end{solucion}

\begin{solucion}{eje:6-embeddings}
\datos GPT-2 small: $V=50\,257$, $\dmodel=768$, $124\,439\,808$ parámetros. Llama~2 7B: $V=32\,000$, $\dmodel=4096$, unos $6738$~millones de parámetros.
\planteamiento la matriz de embeddings tiene $V\cdot\dmodel$ parámetros; sin pesos atados hay dos matrices de esa forma.
\pasos
\begin{enumerate}[label=\arabic*.]
  \item (a) $50\,257\cdot 768=38\,597\,376$; fracción: $38\,597\,376/124\,439\,808=0{,}3102\approx 31{,}0\,\%$.
  \item (b) $124\,439\,808+38\,597\,376=163\,037\,184$ parámetros.
  \item (c) Una matriz: $32\,000\cdot 4096=131\,072\,000$; dos: $262\,144\,000$.
  \item Porcentaje: $262\,144\,000/(6738\cdot 10^{6})\approx 0{,}0389=3{,}89\,\%$.
\end{enumerate}
\resultado \textbf{(a) $38\,597\,376$ ($31\,\%$); (b) $163\,037\,184$; (c) $262\,144\,000$, alrededor del $3{,}9\,\%$ de Llama~2 7B.}
\comprobacion en un modelo pequeño la tabla de embeddings es una fracción grande del total; al crecer $\dmodel$ y el número de capas, el resto de parámetros crece más deprisa que $V\dmodel$, y la fracción disminuye.
\end{solucion}

\begin{solucion}{eje:6-atencion}
\datos forma $(8,12,1024,1024)$; 2 bytes por elemento.
\planteamiento número de elementos $=$ producto de los tamaños de los ejes.
\pasos
\begin{enumerate}[label=\arabic*.]
  \item $8\cdot 12=96$; $1024\cdot 1024=1\,048\,576$; $96\cdot 1\,048\,576=100\,663\,296$ elementos.
  \item Memoria: $100\,663\,296\cdot 2=201\,326\,592$ bytes $\approx 201{,}3$~MB ($192$~MiB).
  \item Con $T=2048$, el factor $T^{2}$ pasa de $1024^{2}$ a $2048^{2}=4\cdot 1024^{2}$: la memoria se multiplica por~4.
\end{enumerate}
\resultado \textbf{$100\,663\,296$ elementos, unos $201$~MB por capa; duplicar el contexto cuadruplica la memoria.}
\comprobacion $201\,326\,592/2^{20}=192$ exactamente, porque $96\cdot 2^{20}\cdot 2/2^{20}=192$. El crecimiento cuadrático con $T$ es la razón de ser de las implementaciones de atención que no almacenan la matriz completa (módulo~17).
\end{solucion}

\begin{solucion}{eje:6-orden}
\datos $A\in\R^{2^{10}\times 2^{16}}$, $B\in\R^{2^{16}\times 2^{5}}$, $C\in\R^{2^{5}\times 2^{14}}$.
\planteamiento un producto de $m\times p$ por $p\times q$ cuesta $mpq$ multiplicaciones y produce $mq$ elementos.
\pasos
\begin{enumerate}[label=\arabic*.]
  \item $(AB)C$: $AB$ cuesta $2^{10}\cdot 2^{16}\cdot 2^{5}=2^{31}$ y es de $2^{10}\times 2^{5}$ ($2^{15}=32\,768$ elementos). Después, $(AB)C$ cuesta $2^{10}\cdot 2^{5}\cdot 2^{14}=2^{29}$.
  \item Total: $2^{31}+2^{29}=5\cdot 2^{29}=2\,684\,354\,560\approx 2{,}68\cdot 10^{9}$.
  \item $A(BC)$: $BC$ cuesta $2^{16}\cdot 2^{5}\cdot 2^{14}=2^{35}$ y es de $2^{16}\times 2^{14}$ ($2^{30}\approx 1{,}07\cdot 10^{9}$ elementos, 4~GiB en 32~bits). Después, $A(BC)$ cuesta $2^{10}\cdot 2^{16}\cdot 2^{14}=2^{40}$.
  \item Total: $2^{40}+2^{35}=33\cdot 2^{35}\approx 1{,}134\cdot 10^{12}$.
  \item Cociente: $(33\cdot 2^{35})/(5\cdot 2^{29})=33\cdot 64/5=422{,}4$.
\end{enumerate}
\resultado \textbf{conviene $(AB)C$: unas 422 veces menos operaciones y una matriz intermedia de $32\,768$ elementos frente a $2^{30}$.} El motivo es que $B$ estrecha la dimensión intermedia a $2^{5}$: conviene aplicarla primero.
\comprobacion ambos órdenes dan el mismo resultado de $2^{10}\times 2^{14}$, por la asociatividad del producto. Es la misma razón por la que LoRA calcula $B(Ax)$ y no $(BA)x$: con $r\ll d$, el primero cuesta $r(d_{\text{in}}+d_{\text{out}})$ multiplicaciones por token.
\end{solucion}

\begin{solucion}{eje:6-cuello}
\datos $Z=HW\T$ (con sesgo, $Z=HW\T+\mathbf{1}_Nb\T$); $L=Z-\ell\mathbf{1}_V\T$ con $\ell_c=\LSE(z_c)$.
\planteamiento cada término de rango~1 añadido aumenta el rango como mucho en~1, porque $\rango(X+Y)\le\rango(X)+\rango(Y)$ (el espacio columna de $X+Y$ está contenido en $\Col(X)+\Col(Y)$).
\pasos
\begin{enumerate}[label=\arabic*.]
  \item (a) $\rango(HW\T)\le d$ (\cref{prop:rango-producto}); $\mathbf{1}_Nb\T$ y $\ell\mathbf{1}_V\T$ tienen rango $\le 1$. Luego $\rango(L)\le d+1+1=d+2$.
  \item (b) $z_1=W h_1=(0,0)$ y $z_2=Wh_2=(1,0)$. $\LSE(0,0)=\ln 2$ y $\LSE(1,0)=\ln(1+e)\approx 1{,}313$.
  \item $L=\begin{pmatrix}-\ln 2&-\ln 2\\1-\ln(1+e)&-\ln(1+e)\end{pmatrix}\approx\begin{pmatrix}-0{,}6931&-0{,}6931\\-0{,}3133&-1{,}313\end{pmatrix}$.
  \item $\det L=(-\ln 2)(-\ln(1+e))-(-\ln 2)(1-\ln(1+e))=\ln 2\ne 0$: $\rango(L)=2$, mientras que $\rango(Z)=1$.
  \item (c) La matriz verdadera es $L^{\star}=-2\ln 2\,J+\ln 2\,I$, con $J$ la matriz de unos: la diagonal vale $\ln\tfrac12=-\ln 2$ y el resto $\ln\tfrac14=-2\ln 2$.
  \item $J$ tiene autovalores $3$, $0$, $0$; luego $L^{\star}$ tiene autovalores $-6\ln 2+\ln 2=-5\ln 2$ y $\ln 2$ (doble), todos no nulos: $\rango(L^{\star})=3$.
  \item Un modelo sin sesgo con $d=1$ produce matrices de log-probabilidades de rango $\le d+1=2$ (\cref{prop:cuello}), que no pueden ser iguales a $L^{\star}$.
\end{enumerate}
\resultado \textbf{(a) $\rango(L)\le d+2$ con sesgo. (b) $\det L=\ln 2$, rango $2=d+1$. (c) $\rango(L^{\star})=3>2$: esas distribuciones no son representables con $d=1$.} \hfill$\blacksquare$
\comprobacion $\det L^{\star}=(-5\ln 2)(\ln 2)^{2}=-5(\ln 2)^{3}\approx -1{,}665$, que coincide con el valor numérico. Con sesgo, la cota del apartado~(a) sería $3$ y el argumento de~(c) ya no se aplicaría: el sesgo añade un grado de libertad de rango~1.
\end{solucion}
